\section{مجازی‌سازی}

\subsection*{بخش آ – اجرای مستقیم چند سیستم‌عامل بدون لایهٔ مجازی‌سازی}

هر سیستم‌عاملِ مهمان برای کار کردن فرض می‌کند که تنها مالک ماشین است و در حالت هسته اجرا می‌شود. سخت‌افزار در زمان بوت در حالت هسته شروع می‌کند و فقط در این حالت می‌توان دستورهای ممتاز (مثل تنظیم زمان‌سنج، مدیریت وقفه، کنترل \lr{I/O} و پاک‌کردن حافظه) را اجرا کرد. بنابراین یک سیستم‌عامل مهمان از سخت‌افزار این انتظارها را دارد:
\begin{itemize}
  \item \textbf{پردازنده:} اجرای آزاد دستورهای ممتاز، تنظیم زمان‌سنج و کنترل‌گر وقفه، ماسک‌کردن وقفه‌ها، تعویض جدول وقفه (\lr{IDT}) و داشتن \lr{CPU} به‌صورت انحصاری. زمان‌بندی را خودش انجام می‌دهد و هیچ دلیلی ندارد پردازنده را به سیستم‌عامل دیگری بدهد.
  \item \textbf{حافظه:} کل حافظهٔ فیزیکی (از آدرس صفر و پیوسته) را از آنِ خود می‌داند، جدول صفحه و \lr{MMU} را خودش مقداردهی می‌کند و می‌تواند هر آدرس فیزیکی را بخواند یا بنویسد.
  \item \textbf{ورودی/خروجی:} دسترسی مستقیم به ثبات‌های کنترل‌گر دستگاه‌ها، درگاه‌ها، خطوط وقفه و \lr{DMA}؛ و انتظار دارد دستگاه‌ها دقیقاً در همان حالتی باشند که خودش رها کرده است.
\end{itemize}

اگر چند سیستم‌عامل هم‌زمان چنین کنترلی داشته باشند، مشکلات زیر پیش می‌آید:
\begin{itemize}
  \item \textbf{تداخل حافظه:} هر هسته روی حافظهٔ دیگری می‌نویسد؛ یک سیستم‌عامل می‌تواند هستهٔ سیستم‌عامل دیگر را پاک کند (همان سناریوی «برنامهٔ کاربر می‌تواند سیستم‌عامل را پاک کند»، این بار میان سیستم‌عامل‌ها).
  \item \textbf{تداخل در \lr{I/O}:} چند سیستم‌عامل هم‌زمان در یک دستگاه می‌نویسند و دستورهای متناقض به یک کنترل‌گر می‌رسد؛ معلوم هم نیست یک وقفه باید به کدام سیستم‌عامل تحویل شود.
  \item \textbf{تداخل در زمان‌بندی:} زمان‌سنج و کنترل‌گر وقفه یکی است و هر سیستم‌عامل آن را به شکل خودش تنظیم می‌کند؛ تنظیمِ دومی تنظیمِ اولی را خراب می‌کند. هیچ‌کس نمی‌تواند \lr{CPU} را از سیستم‌عاملی که وقفه‌ها را غیرفعال کرده یا در حلقهٔ بی‌نهایت است پس بگیرد؛ تقسیم عادلانهٔ منابع ناممکن است.
  \item \textbf{اجرای نادرست:} همه انتظار دارند توسط میان‌افزار بوت شوند و جدول وقفه، جدول صفحه و ثبات‌های کنترلی را در اختیار بگیرند؛ در هر لحظه فقط یک مجموعه از این وضعیت‌ها روی سخت‌افزار می‌تواند فعال باشد و تعویض میان سیستم‌عامل‌ها (ذخیره و بازیابی کل وضعیت ماشین) کار هیچ‌کدام نیست.
  \item \textbf{امنیت و پایداری:} هیچ مرز حفاظتی وجود ندارد؛ خطا یا بدافزار در یک سیستم‌عامل کل ماشین را از کار می‌اندازد یا اطلاعات بقیه را می‌خواند.
\end{itemize}

\textbf{چرا لایهٔ واسط لازم است؟} چون باید یک مرجع واحد وجود داشته باشد که کنترل واقعی منابع را در دست دارد و بقیه را به رعایت قواعد مجبور می‌کند. این لایه همان ناظر ماشین مجازی (\lr{Hypervisor/VMM}) است؛ نرم‌افزار (یا سخت‌افزار)ی است که ماشین‌های مجازی را می‌سازد و اجرا می‌کند (نوع ۱ مستقیماً روی سخت‌افزار) و برای هر مهمان این توهم را می‌سازد که کامپیوتر خصوصی خودش را دارد. وظایف جداسازی آن:
\begin{itemize}
  \item \textbf{پردازنده:} زمان‌بندی \lr{vCPU}ها روی هسته‌های واقعی با زمان‌سنجی که فقط ناظر کنترل می‌کند (پیش‌دستی اجباری؛ به همان دلیل که سیستم‌عامل با زمان‌سنج از انحصار پردازه‌ها جلوگیری می‌کند).
  \item \textbf{حافظه:} نگاشت «فیزیکی مهمان $\rightarrow$ فیزیکی واقعی» با جدول صفحهٔ سایه یا جدول صفحهٔ تودرتو، تخصیص و سهمیه‌بندی حافظه و جلوگیری از دسترسی به حافظهٔ دیگران.
  \item \textbf{ورودی/خروجی:} ارائهٔ دستگاه‌های مجازی (شبیه‌سازی‌شده یا نیمه‌مجازی) یا تخصیص اختصاصی دستگاه همراه با \lr{IOMMU}، و تحویل وقفه‌های مجازی به مهمان درست.
  \item \textbf{حفاظت از خود و مهار خطا:} دستورهای ممتاز فقط در دست ناظر است؛ خرابی یا آلودگی یک مهمان فقط خودش را از کار می‌اندازد و میزبان و مهمان‌های دیگر محافظت می‌شوند.
\end{itemize}

\subsection*{بخش ب – اجرای هستهٔ مهمان مستقیماً در حالت هستهٔ واقعی}

اگر ناظر اجازه دهد هستهٔ مهمان در حالت هستهٔ واقعی اجرا شود، عملاً ناظر کنترل ماشین را از دست داده است:
\begin{itemize}
  \item \textbf{امنیت و دسترسی غیرمجاز:} مهمان می‌تواند هر دستور ممتازی را اجرا کند؛ پس می‌تواند هر آدرس فیزیکی (حافظهٔ ناظر و مهمان‌های دیگر) را بخواند و بنویسد، جدول صفحه و \lr{CR3} را عوض کند یا با \lr{DMA} حافظهٔ دیگران را بخواند. هیچ سازوکاری نمی‌تواند دسترسی غیرمجاز را به تله بیندازد.
  \item \textbf{جداسازی و اثرگذاری روی دیگران:} مهمان می‌تواند جدول وقفه یا کد ناظر را بازنویسی کند و حتی ناظر را دور بزند. چون وقفه‌ها را ماسک و زمان‌سنج را دوباره‌تنظیم می‌کند، ناظر هرگز کنترل \lr{CPU} را پس نمی‌گیرد؛ مهمانِ دارای حلقهٔ بی‌نهایت بقیه را گرسنه نگه می‌دارد (همان مسئلهٔ «\lr{process hogging}» که زمان‌سنج برای جلوگیری از آن وجود دارد).
  \item \textbf{کنترل و تخصیص منابع:} ناظر نمی‌تواند سهمیه بگذارد یا محدودیتی اعمال کند، چون مهمان خود را مالک تمام حافظه و تمام \lr{CPU}ها می‌داند. دستگاه‌های \lr{I/O} به‌طور مستقیم و هم‌زمان توسط چند مهمان برنامه‌ریزی می‌شوند و وقفه‌ها و پیکربندی‌ها با هم تداخل می‌کنند؛ اشتراک‌گذاری کنترل‌شدهٔ دستگاه ناممکن است.
  \item \textbf{پایداری کل سیستم:} خطا در هستهٔ مهمان (اشارهگر خراب، دستور \lr{HLT}، خطای سه‌گانه، پیکربندی غلط سخت‌افزار) مستقیماً ماشین واقعی را متوقف یا بازنشانی می‌کند و فقط به همان مهمان محدود نمی‌ماند. ناظر هم نمی‌تواند وضعیت مهمان را ببیند یا ذخیره کند، پس فعالیت‌هایی مثل عکس‌گیری (\lr{snapshot}) و انتقال زندهٔ ماشین مجازی هم ممکن نیست.
\end{itemize}

نتیجه: هستهٔ مهمان باید از حالت هستهٔ واقعی «تنزل» پیدا کند (مثلاً به حالت کاربر، مانند بخش بعد) یا سخت‌افزار حالت جداگانه‌ای برای مهمان فراهم کند (بخش ه).

\subsection*{بخش ج – روش «تله و شبیه‌سازی» (\lr{Trap-and-Emulate})}

\textbf{روش کار.} ناظر در حالت هستهٔ واقعی اجرا می‌شود و هستهٔ مهمان را در حالت کاربر واقعی اجرا می‌کند؛ مهمان فقط یک «حالت کاربر مجازی» و «حالت هستهٔ مجازی» دارد که هر دو روی حالت کاربر واقعی اجرا می‌شوند. مراحل:
\begin{itemize}
  \item دستورهای معمولی مهمان، چه کد کاربر و چه کد هستهٔ مهمان، مستقیماً و با سرعت تقریباً کامل روی \lr{CPU} اجرا می‌شوند.
  \item وقتی مهمان دستور ممتازی اجرا می‌کند (مثلاً تغییر جدول صفحه، غیرفعال‌کردن وقفه، دستور \lr{I/O})، چون \lr{CPU} در حالت کاربر است، استثنا (تله) رخ می‌دهد و کنترل با ذخیرهٔ وضعیت به ناظر منتقل می‌شود.
  \item ناظر دستور را رمزگشایی می‌کند و اثر آن را روی وضعیت مجازی مهمان (ثبات‌ها و بیت حالت مجازی، جدول صفحهٔ مجازی، دستگاه مجازی) شبیه‌سازی می‌کند و نه روی سخت‌افزار واقعی.
  \item ناظر وضعیت مهمان را بازیابی می‌کند و اجرای مهمان را از دستور بعدی ادامه می‌دهد.
\end{itemize}

\textbf{محدودیت‌ها در معماری‌هایی که برای مجازی‌سازی طراحی نشده‌اند.}
\begin{itemize}
  \item \textbf{دستورهای حساسِ غیرممتاز:} در برخی معماری‌ها (مثل \lr{x86} پیش از افزوده‌شدن پشتیبانی سخت‌افزاری) دستورهایی هستند که وضعیت یا پیکربندی سیستم را می‌خوانند یا تغییر می‌دهند ولی ممتاز نیستند؛ مثل \code{POPF} و \code{SGDT}. این دستورها در حالت کاربر تله ایجاد نمی‌کنند، پس ناظر هرگز مطلع نمی‌شود و نتیجهٔ نادرست به مهمان می‌رسد (تفصیل در بخش د).
  \item مجبور می‌شویم از راه‌حل‌های سنگین‌تر مانند بازنویسی دودویی (ترجمهٔ کد مهمان) یا نیمه‌مجازی‌سازی (تغییر دادن هستهٔ مهمان) استفاده کنیم.
  \item \textbf{هزینهٔ زیاد تله‌ها:} هستهٔ مهمان تعداد زیادی دستور ممتاز اجرا می‌کند و هر کدام به تله می‌افتد؛ این هم باعث افت کارایی می‌شود.
\end{itemize}

\textbf{چرا انتقال‌های زیاد کارایی را کم می‌کند؟} هر انتقال شامل ایجاد استثنا، ذخیرهٔ وضعیت مهمان، سوئیچ به زمینهٔ ناظر، رمزگشایی و شبیه‌سازی دستور، بازیابی وضعیت و بازگشت است. اگر دستور ممتاز روی ماشین واقعی حدود یک چرخهٔ ساعت طول می‌کشد، همین دستور اینجا صدها تا هزاران چرخه هزینه دارد. افزون بر آن، هر انتقال خط لولهٔ \lr{CPU} را خالی می‌کند و کش و \lr{TLB} را آلوده می‌کند. چون هستهٔ مهمان مدام جدول صفحه را به‌روز می‌کند، وقفه‌ها را فعال و غیرفعال می‌کند و به ثبات‌های دستگاه دسترسی می‌دهد (و هر عملیات \lr{I/O} چندین دسترسی به ثبات دارد)، این هزینه چندین برابر می‌شود.

\subsection*{بخش د – ناتوانی معماری در تله‌انداختن همهٔ رفتارهای حساس}

بر اساس شرط پوپک–گلدبرگ، معماری‌ای برای مجازی‌سازی با تله و شبیه‌سازی مناسب است که مجموعهٔ دستورهای حساس (آن‌هایی که وضعیت و پیکربندی منابع را می‌خوانند یا عوض می‌کنند یا بسته به حالت رفتار متفاوت دارند) زیرمجموعهٔ دستورهای ممتاز باشد. اگر چنین نباشد:
\begin{itemize}
  \item \textbf{ایجاد مشکل در پیاده‌سازی:} مهمان دستور حساس را در حالت کاربر مستقیماً اجرا می‌کند و ناظر خبردار نمی‌شود. یا دستور بی‌صدا نتیجهٔ اشتباه می‌دهد (مثلاً \code{POPF} در حالت کاربر تغییر پرچم وقفه را نادیده می‌گیرد، در حالی که مهمان فکر می‌کند وقفه‌ها را غیرفعال کرده است و منطق هسته‌اش خراب می‌شود)، یا اطلاعات واقعی ماشین را به مهمان نشان می‌دهد (مثلاً \code{SGDT} آدرس واقعی جدول توصیف‌گر سراسری را برمی‌گرداند و نه آدرس مجازی). در نتیجه ناظر دیگر نمی‌تواند به مهمان توهم «یک ماشین کامل» بدهد (ویژگی هم‌ارزی).
  \item \textbf{نقض جداسازی:} اگر اثر دستور روی حالت واقعی سخت‌افزار برود، یک مهمان بر وضعیت مشترک ماشین اثر می‌گذارد و بقیهٔ ماشین‌های مجازی و حتی ناظر را تحت تأثیر قرار می‌دهد. اگر دستور اطلاعات بخواند، اطلاعات ناظر یا مهمان‌های دیگر نشت می‌کند و مهمان می‌فهمد روی ماشین مجازی اجرا می‌شود. در هر دو حالت، ناظر کنترل کامل منابع را ندارد (ویژگی کنترل منابع) و مرز حفاظتی شکسته می‌شود.
  \item \textbf{چرا تشخیص صحیح عملیات حساس نیازمندی اصلی است؟} چون هر سه ویژگی مورد نیاز مجازی‌سازی، یعنی وفاداری (مهمان همان رفتار ماشین واقعی را ببیند)، کنترل کامل منابع توسط ناظر و کارایی (بیشتر دستورها مستقیم اجرا شوند)، وقتی برقرار است که ناظر دقیقاً در لحظه‌های لازم دخالت کند. تشخیص ناقص یعنی یا ناظر دخالت نمی‌کند (خطای رفتار و نقض جداسازی)، یا مجبور است همه‌چیز را شبیه‌سازی کند (کارایی ضعیف).
\end{itemize}

\subsection*{بخش ه – پشتیبانی سخت‌افزاری از مجازی‌سازی}

پردازنده‌های جدید (مثل \lr{Intel VT-x} و \lr{AMD-V}) یک حالت اجرای جدید اضافه می‌کنند که «حالت میزبان/ناظر» در برابر «حالت مهمان» نامیده می‌شود. ناظر در حالت میزبان اجرا می‌شود و مهمان در حالت مهمان؛ هستهٔ مهمان در حالت مهمان، حالت هستهٔ مجازی خودش را دارد و می‌تواند بیشتر دستورهای ممتاز را مستقیماً اجرا کند. هر رویدادی که ناظر باید در آن دخالت کند (مثلاً دسترسی به دستگاه، تغییر بعضی ثبات‌های کنترلی، خطای صفحهٔ تودرتو) با تنظیمات ناظر باعث «خروج از ماشین مجازی» (\lr{VM exit}) می‌شود و \lr{CPU} به‌طور خودکار وضعیت مهمان را در ساختار سخت‌افزاری (\lr{VMCS/VMCB}) ذخیره و وضعیت ناظر را بارگذاری می‌کند.

\textbf{تأثیر بر طراحی و عملکرد ناظر:}
\begin{itemize}
  \item \textbf{ساده‌تر شدن طراحی:} ناظر دیگر لازم نیست کد مهمان را اسکن یا بازنویسی کند و کد هسته‌ی مهمان را هم تغییر نمی‌دهد (بدون بازنویسی دودویی و بدون نیمه‌مجازی‌سازی). ناظر به یک حلقهٔ رویدادمحور تبدیل می‌شود: ورود به مهمان، خروج با علت مشخص، پردازش علت (شبیه‌سازی دستگاه یا مدیریت منابع) و ورود دوباره.
  \item چون همهٔ رفتارهای حساس دیگر قابل‌تله‌انداختن‌اند، مشکل دستورهای حساسِ غیرممتاز از بین می‌رود.
  \item ذخیره و بازیابی وضعیت مهمان را سخت‌افزار انجام می‌دهد و ناظر مجبور نیست آن را دستی ساخته و نگه دارد.
  \item جدول صفحهٔ تودرتو (\lr{EPT/NPT}) ترجمهٔ آدرس مهمان را به سخت‌افزار می‌سپارد و ناظر از نگهداری جدول صفحهٔ سایه خلاص می‌شود.
  \item \textbf{کارایی:} مهمان‌های بدون تغییر با سرعت نزدیک به بومی اجرا می‌شوند و فقط رویدادهای واقعاً حساس هزینهٔ خروج را می‌پردازند.
\end{itemize}

